% !TEX root = main.tex
\section{Términos L-S}\label{sec: terminos L-S}
Supongamos que tenemos $N$ electrones en un mismo subnivel $l$, y queremos obtener los términos L-S correspondientes. Es decir, partimos de una configuración electrónica de la forma $l^{N}$ donde asumimos que $N \leq (1/2)(2l + 1)\times (2s+1)$, pues por la simetría ``electrón-hueco'' el caso de $N > (1/2)(2l + 1)\times (2s+1)$ se puede obtener a partir del caso de $N' = (2l + 1)\times (2s+1) - N$ electrones. Por teoría del momento angular, sabemos que el número cuántico total de momento orbital $L$ puede tomar valores entre $L= N \times l$ y $ L =0$ ó $L=l$ según $N$ sea par o impar respectivamente. Esto nos da un primer acercamiento de los valores de $L$ que podemos esperar en los términos L-S. \par
Una parte fundamental del método es estudiar el spin. Análogamente al momento angular, el spin total tomará valores entre $S = N/2$ y $S = 0$ ó $S = 1/2$ según $N$ sea par o impar. Como las posibilidades de la tercera componente del spin individual son menores a las del momento orbital (individual) entonces le asignamos a este un rol secundario. De esta forma tenemos el siguiente panorama: para cada combinación $(M_{L},M_{S})$ tenemos un conjunto de posibles <<tuplas>>\footnote{Nótese que en general no hablamos de una correspondencia entre estados $\ket{M_{L},M_{S}}$ con estados $\ket{m_{1}, \sigma_{1};m_{2},\sigma_{2};\dots ; m_{N}, \sigma_{N}}$, pues en general $\ket{M_{L}, M_{S}}$ será una combinación lineal de estos últimos.}  de los $e^{-}$ individualmente $(m_{1}^{\sigma_{1}}, m_{2}^{\sigma_{2}},\dots m_{N}^{\sigma_{N}})$ donde $m_{i}$ es la tercera componente del momento orbital del electrón $i$ y $\sigma_{i}$ es su tercera componente del spin.  Además, por teoría del momento angular deben cumplir: 
\begin{align}
  \sum_{i=1}^{N} &m_{i} = M_{L} & \sum_{i=1}^{N} &\sigma_{i} = M_{S} \label{eq: suma de 3as componentes} 
\end{align}
Cada uno de estos $m_{i}^{\sigma_{i}}$ son distintos para cumplir el principio de exclusión de Pauli. Existen $2\times(2l + 1))$ posibilidades de $m_{i}^{\sigma_{i}}$ para cada electrón, y por lo tanto, el número de formas posibles de tomar $N$ electrones distintos de entre estas $2\times(2l + 1))$ posibilidades es:
\begin{align}
\# \text{combinaciones} = \binom{2 \times (2l+1)}{N}\label{eq: numero de combinaciones final}
\end{align}
Lo cual se puede comprobar que coincide con el resultado del método tradicional, pues de hecho, es una explicación de dicho resultado. Esto también nos quiere decir que la degeneración acumulada de los términos L-S que obtendremos ha de coincidir con este número de combinaciones, lo cual será una prueba de verificación del método. \par 

A continuación centrémonos en $M_{L}$ y busquemos el número más alto posible que puede tener. Por \eqref{eq: suma de 3as componentes} tendremos que buscar los máximos valores posibles de $m_{i}$ respetando, no obstante, el principio de exclusión. Para ello introducimos la siguiente regla general que nos será útil para el resto del método:\par 
\begin{center} 
\fbox{
  \begin{minipage}[t]{0.9\textwidth}
  Denominando a la expresión $M_{L} = m_{1} + m_{2} + \dots + m_{N}$ como una \textbf{partición} de $M_{L}$, entonces, en ella no podemos repetir el mismo $m$ 2 veces.
  \end{minipage}
}
\end{center}
\vspace*{0.2cm} 
La partición obtenida para el máximo valor de $M_{L}$ se corresponderá con el máximo valor de $L$. \par 
Así, nos queda averiguar el spin. Para ello, deberes sumar las $\sigma$'s de las cuales, por supuesto, no hemos llevado su contabilidad, porque no es necesario. Para la partición de $M_{L}$ más alto, es sencillo averiguar $M_{S}$ porque tendremos muchas repeticiones de $m$'s lo cual significa que las terceras componentes de spin se anulan entre sí y solo tendremos que fijarnos en la $m$ no repetida para averiguar $M_{S}$. Es decir, renombrando los valores que toma $\sigma$ como $+1/2 \to  +$ $-1/2 \to  -$ tendremos una forma tal que: 
\begin{align*}
  M_{L} &= \underbrace{a^{+} + a^{-} + b^{+} + b^{-} + \dots (+ m^{+})}_\text{$N$ sumandos}\; , \quad \text{ donde } a=\max(m_{i})>b>\dots > m 
\end{align*}
Nótese que asignamos los $+$ antes que los $-$, lo cual es una convención, pero no nos hace falta considerar el signo opuesto porque el resultado será simétrico y será irrelevante para obtener $S$. \par
Para $N$ par, no existirá esa $m$ no repetida. Entonces $M_{S} = 0$, lo cual inducirá a que $S=0$ para ese término de máxima $L$. Y para el caso de $N$ impar, existirá ese término de $m$ no repetida y entonces $M_{S} = 1/2$, lo cual inducirá a que $S=1/2$ para ese término de máxima $L$. De ahí podremos extraer el término L-S correspondiente con su degeneración, que será importante para la contabilidad de los siguientes términos. \par
Aunque en general se puede deducir para cada ejemplo, como nota de interés, el lector puede comprobar que en general el valor máximo de $L$, dado  $N$ y $l$ se puede obtener teniendo en cuenta que un valor dado de $m_{l}$ puede como máximo repetirse dos veces. De este modo: 
\begin{align}
   L_\text{max} =\underbrace{2l}_{1} + \underbrace{2(l-1)}_{2}+\ldots+\underbrace{2\left(l-\text{floor}(\frac{N}{2})-1 \right)}_{\text{floor}\left(\frac{N}{2}\right) }+\left[l-\text{floor}\left(\frac{N}{2}\right)-1-1\right]\times \text{mod}(N,2)\\
   L_\text{max} = 2l\times \text{floor}\left(\frac{N}{2}\right)-2\frac{\left(\text{floor}\left(\frac{N}{2}\right)+1\right)\left(\text{floor}\left(\frac{N}{2}\right)+2\right)}{2}+\left[l-\text{floor}\left(\frac{N}{2}\right)-2\right]\times \text{mod}(N,2)   
 \end{align}

Para obtener los siguientes términos, formalmente lo siguiente es aplicar el operador $L_{-}$ para reducir en 1 el valor de $M_{L}$. Si bien recordamos, dicho operador viene dado por: 
\begin{align*}
L_{-} = \sum_{i=1}^{N} l_{-}^{(i)}
\end{align*}
Y si aplicamos $L_{-}$ a un estado $\ket{m_{1},\sigma_{1}; m_{2},\sigma_{2}; \dots m_{N},\sigma_{N}}$ con $M_{L} = \sum_{i=1}^{N} m_{i}$ definido (aunque no tenga un $L$ definido), entonces el resultado será un estado con $M_{L} - 1$. Pero además: 
\begin{align*}
L_{-} \ket{m_{1},\sigma_{1}; m_{2},\sigma_{2}; \dots m_{N},\sigma_{N}} = \sum_{i=0}^{N} l_{-}^{(i)} \ket{m_{1},\sigma_{1}; m_{2},\sigma_{2}; \dots m_{N},\sigma_{N}} = \sum_{i=0}^{N} a_{i} \ket{\dots m_{i} - 1, \sigma_{i} \dots}
\end{align*}
Donde $a_{i}$ por construcción será tal que cumpla el principio de exclusión de Pauli, es decir, $a_{i} = 0$ si $m_{i} - 1$ ya está repetido. \par 
En el lenguaje simple de particiones esto se traduce en que a cada sumando de la partición anterior podemos reducirlo en una cifra para obtener una partición de $M_{L}-1$, siempre y cuando esa reducción no implique repetir más de 2 veces un sumando. Esto dará lugar a varias particiones y tras varias reiteraciones, es posible que a la hora de seguir aplicando $L_{-}$ se obtenga una partición repetida. En ese caso, no la contamos 2 veces, simplemente la ignoramos. Así podremos generar todas las particiones de $M_{L}-1$ a partir de las particiones de $M_{L}$, y sistemáticamente llegar a $M_{L} = 0$. \par 
Para cada uno de los $M_{L}$ tendremos que buscar los posibles $M_{S}$ correspondientes. En ese caso tendremos que generalizar lo que vimos para el caso de $M_{L}$ máximo. En estas particiones tendremos sumandos que se repiten y otros que no. De nuevo, para obtener $M_{S}$ solo nos interesa los términos no repetidos, puesto que los repetidos anulan la tercera de componente de spin entre sí. Asumamos que existen $k$ sumandos no repetidos y sabemos que el número de $M_{S}>0$ distintos será $\mathrm{floor}\qty(k/2) + 1$ porque cada sumando no repetido puede ser $+$ o $-$. Entonces, para una cierta partición $j$, nos interesa conocer el número de combinaciones distintas de $(+)$ y $(-)$ que podemos asignar a los $k$ sumandos no repetidos para un cierto $M_{S}>0$ dado, pues será un indicativo de a qué término $L-S$ corresponde. Para conocer ese número de combinaciones o degeneración, $G^{(j)}_{k}(M_{S})$, simplemente aplicamos combinatoria: 
\begin{align}
  G^{(j)}_{k}(M_{S}) &= \binom{k}{k/2 - M_{S}} \text{ con } k/2 \geq M_{S} > 0 \label{degeneracion M_S}\\ 
  \text{ Por completitud: } G^{(j)}_{k}(M_{S}) &= 0 \text{ para } M_{S} > k/2
\end{align}
Que se puede entender de la siguiente forma: podemos centrarnos en los $(-)$ que debe haber para dar un cierto $M_{S}>0$, entonces, podemos pensar en este resultado como las distintas formas que hay de asignar los $(-)$ a los $k$ sumandos no repetidos. Debemos tener en cuenta que $0! = 1$ es decir que $G_{k}(k/2) = 1$.\par 
Así, para cada partición de $M_{L}$ podremos tener una o varias posibilidades de $M_{S}$. Y esto es esperable, puesto que no solo un cierto $M_{L}$ puede tener varios $L$ asociados, los cuales pueden tener un $S$ distinto; sino que también podemos tener para un mismo $L$ con distintos $S$ asociados. El conteo de particiones de $M_{L}$ con sus correspondientes degeneraciones $G_{k}(M_{S})$ nos ha de dar como mínimo la degeneración total de los términos L-S obtenidos en las particiones de $M_{L}$ anterior y posiblemente, un excedente que se corresponden con términos L-S nuevos de $L$ menor.\par 
Debe tenerse en cuenta que se pueden encontrar términos L-S nuevos repetidos. No es ningún error, es simplemente que hay varias configuraciones de los electrones que pueden dar el mismo término L-S. Estos términos los marcaremos como $K \times L^{S}$ donde $K$ es el número de veces que se ha encontrado ese término L-S. \par 
Merece la pena introducir el concepto paralelo de \textit{número de duplas iguales} $\mathcal{D}$, a las que nos referiremos simplemente como \textit{duplas}. Estas no son más que el número de parejas iguales en una cierta partición de $M_{L}$. Se encuentran relacionadas con el número de sumandos no repetidos $k$ de la siguiente forma: 
\begin{align*}
N = 2 \mathcal{D} + k 
\end{align*}
Donde $N$ es el número de electrones. En los ejemplos que comentemos en los puntos siguientes emplearemos las duplas para establecer \textbf{reglas} que faciliten la contabilidad y obtención de los términos L-S.\par
Finalmente, el método se da por terminado cuando llegamos a $M_{L} = 0$. En ese caso, la degeneración total del número de términos L-S encontrados debe coincidir con el número de combinaciones \eqref{eq: numero de combinaciones final}, lo cual es una prueba de que el método se ha aplicado correctamente. \par


